\documentclass{article}
\usepackage[utf8]{inputenc}
\usepackage{enumitem}
\title{Designvalg og begrunnelser}
\author{}
\date{}

\begin{document}
\maketitle

\section{Overordnet topologi}
Vi har valgt en enkel kondensatorbank (én stor kondensator) fremfor en Marx-generator av følgende grunner:
\begin{itemize}
\item \textbf{Energi:} En Marx-generator med små kondensatorer (f.eks. 100 nF) lagrer svært lite energi per trinn. For å nå 80 J måtte vi hatt svært mange trinn eller store kondensatorer, noe som øker kompleksiteten.
\item \textbf{Strøm:} En enkelt stor kondensator gir en lavere spenning (400 V vs. flere kV), men til gjengjeld en enorm strøm (9 kA). Magnetfeltet avhenger av strømmen, ikke spenningen, så dette er gunstig.
\item \textbf{Enkelhet:} Færre komponenter, færre feilkilder. V2-HIGHPOWER-prosjektet viste at synkronisering av mange gnistgap og høye effekttap i lademotstander kan være problematisk.
\item \textbf{Håndholdbarhet:} En enkelt 1000 µF kondensator er kompakt nok for et håndholdt kabinett.
\end{itemize}

\section{Valg av ladekrets}
En selvoscillerende flyback-omformer er valgt fordi den:
\begin{itemize}
\item Er enkel og krever ingen mikrokontroller eller PWM-kontroller.
\item Kan lade kondensatoren fra et lavspent batteri (54 V) til 400 V.
\item Bruker få komponenter (MOSFET, transformator, diode, zener, motstander).
\end{itemize}
Ulempen er at den ikke har regulering; ladespenningen bestemmes av batterispenningen og transformatorens viklingsforhold. Vi overvåker spenningen manuelt med en LED og stopper ladingen ved å slippe knappen.

\section{Trigger-mekanisme}
Et trigget gnistgap med piezotennmekanisme er valgt fordi:
\begin{itemize}
\item Det gir full kontroll over avfyringstidspunktet.
\item Piezotennere er ekstremt pålitelige og genererer >10 kV uten ekstern strømforsyning.
\item Ingen ekstra elektronikk er nødvendig.
\end{itemize}

\section{Spoledesign}
En flat spiralspole (pancake coil) er valgt fordi den:
\begin{itemize}
\item Gir et konsentrert magnetfelt foran spolen.
\item Har lav induktans, noe som gir høy $di/dt$ og dermed høy indusert spenning.
\item Er enkel å lage med 3D-printet form.
\end{itemize}
Antall vindinger (6) er et kompromiss: færre vindinger gir lavere induktans og høyere strøm, men også lavere magnetfelt for samme strøm. 6 vindinger gir et felt på ~0.7 T, som er mer enn nok.

\section{Komponentvalg}
\begin{itemize}
\item \textbf{Kondensator:} 1000 µF, 450 V, fotoflash-type. Disse er spesielt designet for høye pulsstrømmer og har lav ESR.
\item \textbf{MOSFET:} IRF840 (500 V, 8 A) er en robust, lett tilgjengelig MOSFET.
\item \textbf{Diode:} UF4007 er en rask diode med 1000 V sperrespenning.
\item \textbf{Gnistgap:} Wolframelektroder fra TIG-sveising er ideelle på grunn av høy smeltetemperatur og god erosjonsmotstand.
\item \textbf{Spole:} 10 AWG emaljert kobbertråd gir lav motstand og høy strømtåleevne.
\end{itemize}

\section{Sikkerhet}
\begin{itemize}
\item \textbf{Manuell utladning:} En trykknapp med 1 kΩ / 50 W motstand lar brukeren lade ut kondensatoren trygt på ~5 sekunder.
\item \textbf{Ingen passiv utladning:} For å forenkle kretsen er det ingen bleeder-motstander. Dette betyr at kondensatoren kan holde spenningen i timevis etter lading. Brukeren må alltid sjekke spenningen og lade ut manuelt før berøring.
\item \textbf{LED-indikator:} En grønn LED med 220 kΩ / 1 W motstand lyser når spenningen overstiger ca. 200 V, og gir en visuell advarsel. **Merk: LED-en er ikke en pålitelig indikator for utladet tilstand.**
\item \textbf{Sikring:} En 2 A sikring i batterikretsen beskytter mot kortslutning. **Sikringen beskytter ikke mot lagret energi i kondensatoren.**
\item \textbf{Gatebeskyttelse:} En 1N4148 diode i serie med zeneren hindrer negativ gatespenning.
\end{itemize}

\section{Testversjon}
Før bygging av hovedversjonen anbefales det å bygge en liten 9 V demonstrator (se \texttt{test\_version/}). Denne validerer RLC-utladningsprinsippet og gir praktisk erfaring med spoleinduksjon, uten høyspenningsfare.
\end{document}
